\subsection{Fibonacci--Lie 共振梯：把规范群结构解释为分辨率窗口的离散共振}\label{subsec:fib-lie-resonance}

\paragraph{稳定类型维数的“非外加”离散锚点}
在黄金均值语法（禁止相邻 \(11\)）下，长度 \(m\) 的稳定类型集合满足
\(
|X_m|=F_{m+2}
\)。
因此当我们把“规范自由度的生成元数/玻色子数”视为某一层可被读出的离散维数接口时，
一个自然且可审计的候选是：在若干小窗口上，\(|X_m|\) 与紧李代数的维数发生精确相等，
从而把 \(\mathrm{SU}(2)\)、\(\mathrm{SU}(3)\) 等结构解释为分辨率增长过程中的\emph{共振点}，而非外加假设。

\paragraph{严格结论：Fibonacci--Lie 共振梯（可复现穷举）}
表 \ref{tab:fibonacci_lie_resonance_ladder} 列出在 \(m\le 500\) 的穷举扫描中出现的全部精确共振点：
\begin{table}[H]
\centering
\scriptsize
\setlength{\tabcolsep}{6pt}
\caption{Fibonacci--Lie 共振梯：稳定类型维数 $|X_m|=F_{m+2}$ 在小窗口上出现对紧李代数维数的精确共振。本表列出在 $m\in[1,500]$ 的穷举扫描里出现的全部精确共振点（其余 $m$ 不满足等号）。}
\label{tab:fibonacci_lie_resonance_ladder}
\begin{tabular}{r r l}
\toprule
$m$ & $|X_m|=F_{m+2}$ & 维数共振 \\
\midrule
2 & 3 & $\dim\mathfrak{su}(2)=3;\ \dim\mathfrak{so}(3)=3;\ \dim\mathfrak{sp}(2)=3$\\
4 & 8 & $\dim\mathfrak{su}(3)=8$\\
6 & 21 & $\dim\mathfrak{so}(7)=21;\ \dim\mathfrak{sp}(6)=21$\\
8 & 55 & $\dim\mathfrak{so}(11)=55;\ \dim\mathfrak{sp}(10)=55$\\
\bottomrule
\end{tabular}
\end{table}

并且更强地，若把等号直接写成二次丢番图方程，则在同一穷举范围内只出现极少数解：
% AUTO-GENERATED by scripts/exp_fibonacci_lie_resonance_ladder.py
\[
\begin{aligned}
\{(m,N):1\le m\le 500,\ F_{m+2}=\tfrac{N(N-1)}{2}\} &= \{(2,3),(6,7),(8,11)\},\\
\{(m,N):1\le m\le 500,\ F_{m+2}=N^2-1\} &= \{(2,2),(4,3)\}.
\end{aligned}
\]

上述两段输出均由脚本 \texttt{scripts/exp\_fibonacci\_lie\_resonance\_ladder.py} 生成，
属于可审计、可复现的算术事实（不依赖任何解释层叙事）。

\begin{corollary}[Fibonacci--Lie 共振点的稀缺性与 \texorpdfstring{$\mathfrak{su}(2),\mathfrak{su}(3)$}{su2,su3} 的算术极小性]\label{cor:fib-lie-resonance-scarcity-su2-su3}
在 \(m\le 500\) 的穷举扫描范围内，满足
\(
|X_m|=F_{m+2}=\dim(\mathfrak g)
\)
且 \(\mathfrak g\) 为紧简单 Lie 代数（或等维对偶型）的共振点仅出现在
\[
m\in\{2,4,6,8\},
\qquad
F_{m+2}\in\{3,8,21,55\}
\]
（表 \ref{tab:fibonacci_lie_resonance_ladder}）。
其中 \(3\) 与 \(8\) 在紧简单分类中分别唯一对应 \(\mathfrak{su}(2)\) 与 \(\mathfrak{su}(3)\)；
因此在该离散算术接口下，\(\mathfrak{su}(2)\oplus\mathfrak{su}(3)\) 作为非阿贝尔可见部分体现为两次最早出现的非平凡共振。
\leanverified{paper\_fib\_lie\_resonance\_scarcity\_su2\_su3}
\end{corollary}
\begin{proof}
由表 \ref{tab:fibonacci_lie_resonance_ladder} 直接得到共振点列表；而 \(\dim=3,8\) 的紧简单同构类唯一性为标准分类事实。证毕。
\end{proof}

上述穷举表给出的是有限窗口证书；下面说明，这个有限列表实际上已经是全局完备的。

\begin{theorem}[\texorpdfstring{$A$}{A} 型共振的全局完备性]\label{thm:fib-lie-resonance-global-a-type}
设 \(m\ge 1\)、\(N\ge 2\)。
若
\[
|X_m|=F_{m+2}=\dim\mathfrak{su}(N)=N^2-1,
\]
则唯一可能为
\[
(m,N)=(2,2)
\qquad\text{或}\qquad
(m,N)=(4,3).
\]
等价地，\(A\) 型共振在全部分辨率窗口中只出现于
\[
|X_2|=3=\dim\mathfrak{su}(2),
\qquad
|X_4|=8=\dim\mathfrak{su}(3).
\]
\leanverified{fib\_lie\_resonance\_complete}
\end{theorem}
\begin{proof}
由本小节的基本计数律，
\(
|X_m|=F_{m+2}
\)。
若
\(
F_{m+2}=N^2-1
\)，
则
\[
F_{m+2}+1=N^2
\]
是完全平方。
由 Bugeaud--Luca--Mignotte--Siksek 关于 Fibonacci 数与完全幂偏差的分类 \cite{BugeaudLucaMignotteSiksek2008FibonacciPerfectPowerGap}，
若 \(F_n\pm 1\) 是完全幂，则 \(F_n\in\{0,1,2,3,5,8\}\)；
特别地，若 \(F_n+1\) 是平方，则只能有
\[
F_n\in\{0,3,8\}.
\]
这里 \(n=m+2\ge 3\)，故 \(F_n\neq 0\)，于是只剩
\[
F_{m+2}\in\{3,8\}.
\]
解得
\[
F_{m+2}=3 \iff m+2=4 \iff m=2,
\qquad
F_{m+2}=8 \iff m+2=6 \iff m=4.
\]
相应地，
\(
N^2-1=3,8
\)
只给出
\(
N=2,3
\)。
证毕。
\end{proof}

\begin{theorem}[正交/辛型共振的全局完备性]\label{thm:fib-lie-resonance-global-orthogonal-symplectic}
设 \(m\ge 1\)。

若存在 \(N\ge 3\) 使
\[
|X_m|=F_{m+2}=\dim\mathfrak{so}(N)=\frac{N(N-1)}{2},
\]
则唯一可能为
\[
(m,N)=(2,3),\qquad (6,7),\qquad (8,11).
\]

若存在 \(r\ge 1\) 使
\[
|X_m|=F_{m+2}=\dim\mathfrak{sp}(2r)=r(2r+1),
\]
则唯一可能为
\[
(m,r)=(2,1),\qquad (6,3),\qquad (8,5).
\]
\leanverified{paper\_fib\_lie\_resonance\_orthogonal\_symplectic}
\end{theorem}
\begin{proof}
经典维数公式给出
\[
\dim\mathfrak{so}(N)=\frac{N(N-1)}{2}=T_{N-1},
\qquad
\dim\mathfrak{sp}(2r)=r(2r+1)=T_{2r},
\]
其中 \(T_s=s(s+1)/2\) 为第 \(s\) 个三角数。
因此两族维数都落在三角数列中。

另一方面，Ming 证明 Fibonacci 数列中的三角数只有
\[
1,\ 3,\ 21,\ 55
\]
这四个值；可参见原始结果 \cite{Ming1989TriangularFibonacci}，亦可参见回顾 \cite{McDaniel1996TriangularPell}。
由于 \(m\ge 1\) 时 \(F_{m+2}\ge 2\)，故只能有
\[
F_{m+2}\in\{3,21,55\}.
\]
于是
\[
\frac{N(N-1)}{2}\in\{3,21,55\}
\]
只给出
\[
N=3,7,11,
\]
而
\[
r(2r+1)\in\{3,21,55\}
\]
只给出
\[
r=1,3,5.
\]
最后，由
\[
F_4=3,\qquad F_8=21,\qquad F_{10}=55
\]
可得对应窗口恰为
\[
m=2,\ 6,\ 8.
\]
证毕。
\end{proof}

\begin{theorem}[紧简单 Lie 共振的全局分类]\label{thm:fib-lie-resonance-global-classification}
\leanverified{paper\_fib\_lie\_resonance\_global\_classification}
设 \(m\ge 1\)，\(\mathfrak g\) 为紧简单 Lie 代数。
若
\[
|X_m|=F_{m+2}=\dim\mathfrak g,
\]
则参数对 \((m,\mathfrak g)\) 只能属于以下八种情形；反之，这八种情形都确实满足等号：
\[
(2,\mathfrak{su}(2)),\quad
(4,\mathfrak{su}(3)),\quad
(2,\mathfrak{so}(3)),\quad
(6,\mathfrak{so}(7)),\quad
(8,\mathfrak{so}(11)),
\]
\[
(2,\mathfrak{sp}(2)),\quad
(6,\mathfrak{sp}(6)),\quad
(8,\mathfrak{sp}(10)).
\]
按同构类去重后，恰得到六种可能：
\[
\mathfrak{su}(2),\qquad
\mathfrak{su}(3),\qquad
\mathfrak{so}(7),\qquad
\mathfrak{sp}(6),\qquad
\mathfrak{so}(11),\qquad
\mathfrak{sp}(10).
\]
此外，不存在任何 \(D\) 型或例外型共振。
\end{theorem}
\begin{proof}
紧简单 Lie 代数的经典分类只包含四个经典系列 \(A,B,C,D\) 与五个例外型
\[
G_2,\qquad F_4,\qquad E_6,\qquad E_7,\qquad E_8,
\]
其维数分别为
\[
\dim G_2=14,\qquad
\dim F_4=52,\qquad
\dim E_6=78,\qquad
\dim E_7=133,\qquad
\dim E_8=248.
\]

对 \(A\) 型，定理 \ref{thm:fib-lie-resonance-global-a-type} 已穷尽全部可能。
对正交/辛型，定理 \ref{thm:fib-lie-resonance-global-orthogonal-symplectic} 给出
\[
N\in\{3,7,11\},
\qquad
r\in\{1,3,5\}.
\]
因此正交侧除低秩重合 \(\mathfrak{so}(3)\cong\mathfrak{su}(2)\) 外，只出现 \(\mathfrak{so}(7)\) 与 \(\mathfrak{so}(11)\)；
而由于没有任何偶数 \(N\) 出现，故不存在 \(D\) 型共振。
辛侧则得到 \(\mathfrak{sp}(2)\cong\mathfrak{su}(2)\)、\(\mathfrak{sp}(6)\)、\(\mathfrak{sp}(10)\)。

最后只需排除例外型。
注意
\[
F_7=13<14<F_8=21,
\]
\[
F_9=34<52<F_{10}=55,
\qquad
F_{10}=55<78<F_{11}=89,
\]
\[
F_{11}=89<133<F_{12}=144,
\qquad
F_{13}=233<248<F_{14}=377.
\]
故
\[
14,\ 52,\ 78,\ 133,\ 248
\]
都不是 Fibonacci 数，因此不可能等于 \(F_{m+2}\)。
另一方面，前述八种情形分别对应维数
\[
3,\ 8,\ 3,\ 21,\ 55,\ 3,\ 21,\ 55,
\]
而
\[
F_4=3,\qquad F_6=8,\qquad F_8=21,\qquad F_{10}=55,
\]
故它们都确实实现
\(
|X_m|=F_{m+2}=\dim\mathfrak g
\)。
综上，结论成立。证毕。
\end{proof}

\begin{corollary}[表 \ref{tab:fibonacci_lie_resonance_ladder} 的全局闭合与共振梯在 \(m=8\) 处终止]\label{cor:fib-lie-resonance-ladder-global-closure}
表 \ref{tab:fibonacci_lie_resonance_ladder} 在有限扫描窗口 \(1\le m\le 500\) 中观察到的共振点
\[
m\in\{2,4,6,8\}
\]
事实上已经是全体解。
更精确地：
\begin{enumerate}
  \item 若要求 \(A\) 型共振，则只可能在 \(m=2,4\)；
  \item 若要求正交/辛型共振，则只可能在 \(m=2,6,8\)；
  \item 对全部紧简单 Lie 代数而言，不存在任何 \(m>8\) 的共振窗口。
\end{enumerate}
此外，仅有两组相邻偶窗口双共振：
\[
(m,m+2)=(2,4)
\qquad\text{与}\qquad
(m,m+2)=(6,8),
\]
对应维数链分别为
\[
3\to 8,
\qquad
21\to 55.
\]
\leanverified{fib\_double\_gt\_sq\_sub\_one}
\end{corollary}
\begin{proof}
前三条是定理 \ref{thm:fib-lie-resonance-global-a-type}、
\ref{thm:fib-lie-resonance-global-orthogonal-symplectic} 与
\ref{thm:fib-lie-resonance-global-classification} 的直接汇总。
对最后一条，只需从全局分类中读取可共振窗口集合
\[
\{2,4,6,8\}.
\]
其中相差 \(2\) 的仅有 \((2,4)\) 与 \((6,8)\) 两对，而对应维数链即
\[
F_4=3\to F_6=8,
\qquad
F_8=21\to F_{10}=55.
\]
证毕。
\end{proof}
\leanverified{paper\_fib\_lie\_resonance\_ladder\_global\_closure}

\begin{remark}[与边界词塔的四步平移]\label{rem:fib-lie-resonance-shift-by-4}
在第 \ref{subsec:bdry-tower-zeck-gut} 小节中，我们使用端点 \((1,1)\) 纤维的基数
\(
|X_m^{\mathrm{bdry}}|=F_{m-2}
\)
作为“边界模态数”。
注意到
\[
|X_m^{\mathrm{bdry}}|=F_{m-2}=F_{(m-4)+2}=|X_{m-4}|,
\]
因此本小节的共振梯（基于 \(|X_m|\)）在边界塔口径下表现为一个 \(m\mapsto m+4\) 的离散平移：
稳定类型共振点 \(m=2,4,6,8,\dots\) 对应边界层 \(m=6,8,10,12,\dots\)。
这解释了为何边界三层塔（\(m=4,6,8\)）会自然对齐 SM 的 \(1,3,8\) 三段玻色子计数接口：
它是同一 Fibonacci 维数骨架在不同纤维截面上的读出。
\end{remark}

\begin{remark}[最小额外生成元计数：\(21-12=9\)]\label{rem:fib-lie-extra-9}
若把 \(m=6\) 的稳定类型维数 \( |X_6|=21 \) 视为“最小统一层”的生成元数候选，
并把 SM 的可见玻色子数视为 \(12\)（第 \ref{subsec:bdry-tower-zeck-gut} 小节的边界塔接口），
则二者之间的差额为一个刚性的整数：
\[
21-12=9.
\]
这提供一个与“先选大群再分解”的先验群分解路线不同的计数指纹：额外通道必须以 \(9\) 个生成元的离散尺度出现，
其具体实现应通过本框架的扭曲谱/边界层离散标签/投影门机制来审计（而非作为连续可调参数引入）。
\end{remark}

\begin{theorem}[二值手征读出的最小内禀锚点由 window-\(6\) 的 boundary 二层 uplift 强制]\label{thm:window6-chirality-anchor-minimal}
在“分辨率—群论”协议中，把二值手征读出理解为一个\emph{协议内禀}的 \(\ZZ/2\ZZ\)-分级：它必须仅由折叠/uplift 的硬数据给出，不允许引入外置寄存器，并要求与分辨率限制兼容且可由有限证书审计。
则该分级的最小可实现窗口为 \(m=6\)，并且（在同构意义下）唯一等同于 window-\(6\) 的 boundary sheet parity。
\par
因此，Fibonacci--Lie 共振梯在 \(m\in\{2,4\}\) 的 bulk 共振点（表 \ref{tab:fibonacci_lie_resonance_ladder}）只能提供维数接口，
并不能在上述内禀准则下产生非平凡二值手征通道；任何将该手征自由度直接归因于 \(m=2\) 或 \(m=4\) 的 bulk 稳定类型结构的方案，都必须引入额外外置选择，从而不属于该协议口径。
\leanverified{paper\_window6\_chirality\_anchor\_minimal}
\end{theorem}

\begin{proof}
表 \ref{tab:fibonacci_lie_resonance_ladder} 的 \(m=2,4\) 共振仅陈述 \(|X_m|=F_{m+2}\) 与紧李代数维数发生等号，
该等号本身不携带一个由协议硬数据唯一决定的 \(\ZZ/2\ZZ\)-分级。
若要求分级由 uplift 数据内禀恢复，则必须存在无需额外选择的逐纤维二元配对规则。
对 boundary uplift，命题 \ref{prop:xi-window6-sheet-parity-anchor-m6} 给出：当且仅当 \(m=6\) 时，
每条 boundary 纤维为二元集合且非零 uplift 增量唯一（\(\Delta_6^{\mathrm{bdry}}=\{0,34\}\)），从而诱导无歧义的 sheet parity；
而当 \(m=7,8\) 时 boundary 纤维为三元集合且增量集合含两种不同非零元，二元化必引入外置选择，故不满足内禀准则（并见定理 \ref{thm:window6-bdry-sheet-parity-nonfunctorial-extension}）。
因此该二值分级在协议意义下被强制锚定于 \(m=6\) 的 boundary 二层 uplift，且在同构意义下唯一。证毕。
\end{proof}

\begin{proposition}[window-\(6\) 的 dyadic 微观预算三段式分解]\label{prop:window6-dyadic-budget-three-stage}
令 \(\Omega_6:=\{0,1\}^6\) 为 window-\(6\) 的 dyadic 微态空间（\(|\Omega_6|=2^6=64\)），
并令 \(X_6\) 为稳定类型集合（\(|X_6|=F_8=21\)）。
定义 dyadic 缺额
\[
H^{\mathrm{bin}}_6:=|\Omega_6|-|X_6|=64-21=43.
\]
则在 orbit--Levi 骨架给出的分块维数 \(21=15\oplus 3\oplus 3\)（见注 \ref{rem:window6-orbit-levi-threshold-compatible-closure}）与 boundary uplift 主位移 \(\delta=34\)（推论 \ref{cor:terminal-foldbin6-boundary-pure-f9-alias}）的约束下，成立严格恒等式
\[
2^6\ =\ (15+3+3)\ +\ (34+8+1),
\]
等价地
\[
H^{\mathrm{bin}}_6\ =\ \delta+|X^{\mathrm{bdry}}_8|+|X^{\mathrm{bdry}}_4|
\ =\ 34+8+1.
\]
\leanverified{paper\_gut\_window6\_dyadic\_budget\_three\_stage}
\end{proposition}

\begin{proof}
由定义得 \(H^{\mathrm{bin}}_6=64-21=43\)。
推论 \ref{cor:terminal-foldbin6-boundary-pure-f9-alias} 给出 boundary 二层 uplift 的唯一非零位移 \(\delta=34\)。
另一方面，边界塔基数满足 \(|X_m^{\mathrm{bdry}}|=F_{m-2}\)（定理 \ref{thm:boundary-shift4-uplift-isomorphism}），故
\(|X_8^{\mathrm{bdry}}|=F_6=8\) 且 \(|X_4^{\mathrm{bdry}}|=F_2=1\)；
并且标准模型的 Zeckendorf 边界签名为 \(\Sigma_{\mathrm{SM}}=\{4,6,8\}\) 且 \(12=8+3+1\)（表 \ref{tab:zeckendorf_gut_signatures}）。
代入即得
\(
43=34+8+1
\)
以及
\(
64=(15+3+3)+(34+8+1)
\)。
证毕。
\end{proof}

\begin{remark}[fold-gauge 的 \texorpdfstring{$(\ZZ_2)^{21}$}{(Z2)^21} 二值电荷层（\texorpdfstring{$m=6$}{m=6}，dyadic 读出）]\label{rem:fold-gauge-z2-21}
除稳定类型维数 \(|X_6|=21\) 的“维数层”外，\(m=6\) 在纤维冗余层还给出一个独立、可审计的电荷格。
对二进制区间折叠 \(\mathrm{Fold}^{\mathrm{bin}}_6:\{0,\dots,63\}\to X_6\)（定义 \ref{def:fold-bin}），其纤维自同构群
\[
\mathcal{G}^{\mathrm{bin}}_6:=\mathrm{Aut}\!\bigl(\mathrm{Fold}^{\mathrm{bin}}_6\bigr)
\]
满足
\[
\bigl(\mathcal{G}^{\mathrm{bin}}_6\bigr)^{\mathrm{ab}}\cong (\ZZ_2)^{21}
\]
（推论 \ref{cor:fold-bin-gauge-m6}）。
这里每个 \(\ZZ_2\) 因子可由对应纤维内置换的符号表示读出，可解释为“最小非平凡窗口上可见的二值 gauge 电荷”。
该 \(21\) 维二值电荷层与 \S\ref{subsec:bdry-tower-zeck-gut} 的 \(B_3/C_3\) 种子字典在维数上吻合，并为进一步把
\(
21=\dim\wedge^2\RR^7=\dim\mathfrak{so}(7)
\)
这类几何解释升级为可审计约束提供一条额外证据链。
\end{remark}

\begin{remark}[素数 \texorpdfstring{$7,11$}{7,11} 的双重共振对齐：Fibonacci--Lie 共振点 \texorpdfstring{$\leftrightarrow$}{<->} 有限域嵌入模数]\label{rem:fib-lie-primes-7-11-double-alignment}
由推论 \ref{cor:fib-lie-resonance-scarcity-su2-su3}，离散维数接口在小窗口上出现 \(F_{m+2}\in\{21,55\}\) 的两处正交族共振点；
其中 \(21=\dim\mathfrak{so}(7)\) 已在注 \ref{rem:fold-gauge-z2-21} 中显性对齐，而 \(55=\dim\mathfrak{so}(11)\) 为标准维数恒等式。
另一方面，共振窗最小递推特征多项式的有限域谱指纹显示：二阶碰撞核 \(A_2\) 的三次特征式在模 \(7\) 与模 \(11\) 处发生可审计嵌入（命题 \ref{prop:pom-charpoly-modp-a2-embedding}）。
\par
因此 \(p\in\{7,11\}\) 在本文中同时承担两种互不依赖的“共振角色”：一方面作为稳定类型维数与正交李代数维数公式的离散交汇参数（\(21,55\)），另一方面作为碰撞矩递推谱指纹中三次骨架可见化的最小模数。
该双重对齐为后续建立“群论分辨率签名”与“谱指纹嵌入模数”之间的算术桥梁提供一条可扩展且可证伪的起点。
\end{remark}
